% СХЕМА 4 — HMM как цепь состояний
\begin{tikzpicture}[
  every node/.style={font=\footnotesize},
  state/.style={draw, thick, circle, minimum size=0.85cm, inner sep=0pt, fill=white},
  arr/.style={-{Latex[length=2mm]}, thick},
  skip/.style={-{Latex[length=2mm]}, thick, dashed, bend left=45},
]
  \foreach \i/\p in {0/60, 1/62, 2/64, 3/65, 4/67} {
    \node[state] (s\i) at (\i*1.4, 0) {$s_{\i}$};
    \node[font=\scriptsize] at (\i*1.4, -0.7) {$p_{\i}{=}\p$};
  }

  \foreach \i in {0,1,2,3,4} {
    \draw[arr] (s\i.north) to[in=120, out=60, looseness=6]
               node[above, font=\tiny] {stay} (s\i.north);
  }

  \foreach \i [evaluate=\i as \j using int(\i+1)] in {0,1,2,3} {
    \draw[arr] (s\i) -- node[below, font=\tiny] {adv} (s\j);
  }

  \draw[skip] (s0) to node[above, font=\tiny] {skip} (s2);
  \draw[skip] (s2) to node[above, font=\tiny] {skip} (s4);

  \node[font=\footnotesize] (obs) at (3.0, -1.7)
        {observation $o_t \;\sim\; \mathcal{N}(p_i, \sigma^2)$};
  \draw[->, thick] (obs.north) -- (s2.south);
\end{tikzpicture}
